% Part: first-order-logic
% Chapter: syntax-and-semantics
% Section: unique-readability

\documentclass[../../../include/open-logic-section]{subfiles}

\begin{document}

\olfileid{fol}{syn}{unq}

\olsection{एकमेव वाचनीयता}

\begin{explain}
आपण !!{formula}s ची ज्या रीतीने व्याख्या केली आहे त्यामुळे प्रत्येक
!!{formula} चे \emph{एकमेव वाचन} असते; म्हणजे !!{formula}s साठीच्या
आपल्या रचनानियमांनुसार ते तयार करण्याचा मूलतः एकच मार्ग आणि त्याची
``अर्थलावणी'' करण्याचा एकच मार्ग असतो. तसे नसते, तर आपल्याला संदिग्ध
!!{formula}s मिळाली असती, म्हणजे एकाहून अधिक वाचने किंवा अर्थलावण्या
असलेली !!{formula}s---आणि हे आपल्याला टाळायचे आहे हे स्पष्ट आहे. पण
त्याहून महत्त्वाचे म्हणजे, या गुणधर्माविना आपण पुढे देणार असलेल्या
बहुतेक व्याख्या आणि सिद्धता चालणार नाहीत.

एकमेव वाचनीयतेची खात्री न देणारे !!{formula}s तयार करण्याचे सदोष नियम
दिले असते तर काय झाले असते हे पाहणे, हा मुद्दा स्पष्ट करण्याचा बहुधा
सर्वोत्तम मार्ग आहे. उदाहरणार्थ, संयोजकांच्या रचनानियमांत आपण कंस
विसरलो असतो; उदा., पुढील गोष्टीला परवानगी दिली असती:
\begin{quote}
$!A$ आणि $!B$ ही !!{formula}s असतील, तर $!A \lif !B$ हेही तसेच आहे.
\end{quote}
$!D$ या आण्विक सूत्रापासून सुरुवात करून यामुळे $!D \lif !D$ तयार करता
आले असते. यापासून, $!D$ सोबत, $!D \lif !D \lif !D$ मिळाले असते. पण
हे करण्याचे दोन मार्ग आहेत:
\begin{enumerate}
\item $!D$ हे $!A$ आणि $!D \lif !D$ हे $!B$ मानणे.
\item $!A$ हे $!D \lif !D$ आणि $!B$ हे~$!D$ मानणे.
\end{enumerate}
त्याप्रमाणे, !!{formula}~$!D \lif !D \lif !D$ ``वाचण्याचे'' दोन मार्ग
आहेत. ते $!B \lif !C$ या रूपात असून $!B$ हे $!D$ आणि $!C$ हे $!D
\lif !D$ आहे; पण ते $!B \lif !C$ या रूपात \emph{देखील आहे}, जेथे $!B$
हे $!D \lif !D$ आणि $!C$ हे~$!D$ आहे.

असे झाले, तर आपल्या व्याख्या नेहमी चालणार नाहीत. उदाहरणार्थ, एखाद्या
सूत्राचा !!{main operator} परिभाषित करताना आपण म्हणतो: $!B \lif !C$
या रूपातील सूत्रात~$\lif$ ची दर्शवलेली आस्थिति हाच !!{main operator}
आहे. पण $!D \lif !D \lif !D$ या सूत्राची वर सांगितलेल्या दोन वेगळ्या
मार्गांनी $!B \lif !C$ शी जुळणी करता आली, तर एका बाबतीत $\lif$ ची
पहिली आस्थिति !!{main operator} मिळते आणि दुसऱ्या बाबतीत दुसरी
आस्थिति मिळते. पण !!{main operator} हे !!{formula} चे \emph{फलन} असावे
असा आपला उद्देश आहे; म्हणजे प्रत्येक !!{formula} मध्ये !!{main
  operator} ची नेमकी एक आस्थिति असली पाहिजे.
\end{explain}

\begin{lem}
!!a{formula}~$!A$ मधील डाव्या आणि उजव्या कंसांची संख्या समान असते.
\end{lem}

\begin{proof}
$!A$ ज्या रीतीने तयार होते तिच्यावर विगमन करून हे सिद्ध करू. यासाठी
दोन गोष्टी आवश्यक आहेत: (a) प्रथम, सर्व आण्विक सूत्रांना संबंधित
गुणधर्म लागू होतो हे सिद्ध करावे लागते (विगमनाचा आधार). (b) नंतर,
दिलेल्या सूत्रांपासून नवी सूत्रे तयार केली आणि जुन्या सूत्रांना तो
गुणधर्म लागू होत असेल, तर नव्या सूत्रांनाही तो लागू होतो हे सिद्ध
करावे लागते.

$l(!A)$ ही $!A$ मधील डाव्या कंसांची संख्या आणि $r(!A)$ ही उजव्या
कंसांची संख्या असू द्या; त्याचप्रमाणे $l(t)$ आणि $r(t)$ या पद~$t$
मधील अनुक्रमे डाव्या आणि उजव्या कंसांच्या संख्या असू द्या.

\begin{prob}
    कोणत्याही पद~$t$ साठी $l(t) = r(t)$ हे सिद्ध करा.
\end{prob}

\begin{enumerate}
\tagitem{prvFalse}{\indcase{!A}{\lfalse}{$\indfrm$ मध्ये $0$ डावे आणि $0$
    उजवे कंस आहेत.}}{}

\tagitem{prvTrue}{\indcase{!A}{\ltrue}{$\indfrm$ मध्ये $0$ डावे आणि $0$
    उजवे कंस आहेत.}}{}

\item \indcase{!A}{\Atom{R}{t_1,\dots,t_n}}{$l(\indfrm) = 1 + l(t_1) +
  \dots + l(t_n) = 1 + r(t_1) + \dots + r(t_n) = r(\indfrm)$. येथे
  कोणत्याही पद~$t$ साठी $l(t) = r(t)$ ही, सरावासाठी सोडलेली, बाब
  वापरली आहे.}

\item \indcase{!A}{\eq[t_1][t_2]}{$l(\indfrm) = l(t_1) + l(t_2) =
  r(t_1) + r(t_2) = r(\indfrm)$.}

\tagitem{prvNot}{\indcase{!A}{\lnot !B}{विगमन गृहीतकानुसार,
    $l(!B) = r(!B)$. म्हणून $l(\indfrm) = l(!B) = r(!B) =
    r(\indfrm)$.}}{}

\item \indcase{!A}{(!B \ast !C)}{विगमन गृहीतकानुसार, $l(!B) =
  r(!B)$ आणि $l(!C) = r(!C)$. म्हणून $l(\indfrm) = 1 + l(!B) + l(!C) =
  1 + r(!B) + r(!C) = r(\indfrm)$.}

\tagitem{prvAll}{\indcase{!A}{\lforall[x][!B]}{विगमन
    गृहीतकानुसार, $l(!B) = r(!B)$. म्हणून $l(\indfrm) = l(!B) = r(!B) =
    r(\indfrm)$.}}{}

% Print case for \lexists if prvEx and not prvAll
\tagitem{prvAll,notprvEx}{}{\indcase{!A}{\lexists[x][!B]}{विगमन
    गृहीतकानुसार, $l(!B) = r(!B)$. म्हणून $l(\indfrm) = l(!B) = r(!B) =
    r(\indfrm)$.}}

% Just say similarly if prvEx and prvAll
\tagitem{notprvAll,notprvEx}{}{\indcase{!A}{\lexists[x][!B]}{त्याचप्रमाणे.}}
\end{enumerate}
\end{proof}

\begin{defn}[उचित पूर्वखंड]
$!B$ ही चिन्हमाला आणि एक रिकामी नसलेली चिन्हमाला यांची जोडणी केल्यावर
चिन्हमाला~$!A$ मिळत असेल, तर चिन्हमाला $!B$ ही चिन्हमाला~$!A$ चा
\emph{उचित पूर्वखंड} आहे.
\end{defn}

\begin{lem}\ollabel{lem:no-prefix}
$!A$ हे !!a{formula} आणि $!B$ हा $!A$ चा उचित पूर्वखंड असेल, तर $!B$
हे !!a{formula} नाही.
\end{lem}

\begin{proof}
सराव.
\end{proof}

\begin{prob}
\olref[fol][syn][unq]{lem:no-prefix} सिद्ध करा.
\end{prob}

\begin{prop}
\ollabel{prop:unique-atomic}
$!A$ हे आण्विक !!{formula} असेल, तर पुढीलपैकी एक आणि फक्त एकच अट त्याला
लागू होते.
\begin{enumerate}
\tagitem{prvFalse}{$!A \ident \lfalse$.}{}
\tagitem{prvTrue}{$!A \ident \ltrue$.}{}
\item $!A \ident \Atom{R}{t_1,\dots,t_n}$, जेथे $R$ हे $n$-स्थानी
  !!{predicate}, $t_1$, \dots, $t_n$ ही पदे आणि $R$, $t_1$, \dots,
  $t_n$ यांपैकी प्रत्येक एकमेव रीतीने निर्धारित आहे.
\item $!A \ident \eq[t_1][t_2]$, जेथे $t_1$ आणि $t_2$ ही एकमेव रीतीने
  निर्धारित पदे आहेत.
\end{enumerate}
\end{prop}

\begin{proof}
सराव.
\end{proof}

\begin{prob}
\olref[fol][syn][unq]{prop:unique-atomic} सिद्ध करा (सूचना: पदांसाठी
\olref[fol][syn][unq]{lem:no-prefix} ची एक आवृत्ती मांडा आणि सिद्ध करा.)
\end{prob}

\begin{prop}[एकमेव वाचनीयता]
प्रत्येक !!{formula} ला पुढीलपैकी एक आणि फक्त एकच अट लागू होते.
\begin{enumerate}
\item $!A$ हे आण्विक आहे.

\tagitem{prvNot}{$!A$ हे $\lnot !B$ या रूपात आहे.}{}

\tagitem{prvAnd}{$!A$ हे $(!B \land !C)$ या रूपात आहे.}{}

\tagitem{prvOr}{$!A$ हे $(!B \lor !C)$ या रूपात आहे.}{}

\tagitem{prvIf}{$!A$ हे $(!B \lif !C)$ या रूपात आहे.}{}

\tagitem{prvIff}{$!A$ हे $(!B \liff !C)$ या रूपात आहे.}{}

\tagitem{prvAll}{$!A$ हे $\lforall[x][!B]$ या रूपात आहे.}{}

\tagitem{prvEx}{$!A$ हे $\lexists[x][!B]$ या रूपात आहे.}{}
\end{enumerate}
शिवाय, प्रत्येक बाबतीत $!B$, किंवा $!B$ आणि $!C$, एकमेव रीतीने
निर्धारित असतात. म्हणजे, उदाहरणार्थ, $!B$, $!C$ आणि $!B'$, $!C'$ अशा
दोन वेगळ्या जोड्या अस्तित्वात नाहीत की $!A$ हे एकाच वेळी
\iftag{prvIf}{$(!B \lif !C)$ आणि $(!B' \lif !C')$ या दोन्ही रूपांत असेल.}{
    \iftag{prvOr}{$(!B \lor !C)$ आणि $(!B' \lor !C')$ या दोन्ही रूपांत असेल.}{
      \iftag{prvAnd}{$(!B \land !C)$ आणि $(!B' \land !C')$ या दोन्ही रूपांत असेल.}{}}}
\end{prop}

\begin{proof}
रचनानियमांनुसार, !!a{formula} आण्विक नसेल, तर त्याची सुरुवात आरंभीच्या
कंसाने~(, \iftag{prvNot}{$\lnot$ ने,}{} किंवा एखाद्या संख्यापकाने झाली
पाहिजे. दुसरीकडे, पुढीलपैकी एखाद्या चिन्हाने सुरू होणारे प्रत्येक
!!{formula} आण्विक असले पाहिजे: !!a{predicate}, !!a{function},
!!a{constant}\iftag{prvFalse}{, $\lfalse$}{}\iftag{prvTrue}{, $\ltrue$}{}.

म्हणून आपल्याला प्रत्यक्षात एवढेच दाखवायचे आहे: $!A$ हे $(!B \ast
!C)$ या रूपात आणि $(!B' \mathbin{\ast'} !C')$ या रूपातही असेल, तर
$!B \ident !B'$, $!C \ident !C'$ आणि $\ast = {\ast'}$.

म्हणून $!A \ident (!B \ast !C)$ आणि $!A \ident (!B'
\mathbin{\ast'} !C')$ हे दोन्ही गृहीत धरा. मग एकतर $!B \ident !B'$
किंवा तसे नाही. तसे असेल, तर स्पष्टपणे $\ast = {\ast'}$ आणि $!C
\ident !C'$, कारण त्या स्थितीत त्या $!A$ च्या एकाच स्थानी सुरू होणाऱ्या
आणि समान लांबीच्या उपचिन्हमाला आहेत. उरलेली बाब म्हणजे $!B
\not\ident !B'$. $!B$ आणि $!B'$ या दोन्ही $!A$ च्या एकाच स्थानी सुरू
होणाऱ्या उपचिन्हमाला असल्यामुळे, त्यांपैकी एक दुसरीचा उचित पूर्वखंड
असला पाहिजे. पण \olref{lem:no-prefix} नुसार हे अशक्य आहे.
\end{proof}


\end{document}
